\section{Continuum fields at a fixed bare-coupling offset}
\label{sec:fields}

The density comparison of Proposition~\ref{prop:matching} leaves one
real parameter, the limiting offset of the bare coupling.  We now
construct the field hierarchy at each offset and compare the three
blocking prescriptions in common physical coordinates.  The last part
of the section controls the full two-point measure, including its tails.
This is what permits normalization by susceptibility and second-moment
length; convergence on compact tests alone would not suffice.

Throughout, the parameters have the choices of
Propositions~\ref{prop:admission} and~\ref{prop:uniform-volume}.
Write $I=(-R,R)$, and let a run of type $P\in\{r,+,-\}$, depth $N$,
and offset $s$ have bare coupling $\beta=\vartheta_N+s$.
Recall the matching condition
\begin{equation}
 (s_2+A_{P_2})-(s_1+A_{P_1})\longrightarrow0,
 \label{eq:field-offset-matching}
\end{equation}
for two sequences of runs whose depths tend to infinity.  The
discrepancies $u_j$ and $\lambda_j$ are the source-free density and
precise-coupling discrepancies at layer $j$.  By
Proposition~\ref{prop:matching}, both tend to zero at each fixed layer
under \eqref{eq:field-offset-matching}.  All finite-period comparisons
below use the same reference thresholds and the common record lists,
with zero padding, prescribed in Section~\ref{sec:setup}.

\subsection{Local sources and their normalization}

We use the exact source representation of Section~\ref{sec:sources}.
Its additional regular records have positive source degree and are
anchored at a source site; their complete supports include every source
and spin argument.  The covering records retain their mandatory
inventories and complete supports.  All norms sum absolute power-series
coefficients at source radius one.  In particular, identifying source
variables or setting some of them to zero respects the bounds.
The leading linear spin source is kept separately from these records.
Spin-independent terms of source degree at least two stay in their
supported regular records.  There is no source-dependent volume scalar
whose removal could change a moment generating function.

To distinguish a source cap from the offset radius $R$, write
\[
 R_j^{\mathrm{src}}=H_j^{-1/10},\qquad
 w_j=\exp(-p_j^{1/4}).
\]
These are the caps denoted by $R_j,w_j$ in
Equation~\eqref{eq:orig-27}.  In each bulk or finite-period
representation, take the positively weighted source discrepancy of
Proposition~\ref{prop:source-step}, divided by these two caps, with
regular differences measured through seven derivatives and absent
records padded by zero.  Write these quantities as
$v_j^{\mathrm{src,bulk}}$ and $v_j^{\mathrm{src},T}$, where $T$ denotes
a common terminal period followed through the layers.  In parallel
with the density discrepancy \eqref{eq:offset-family-discrepancy}, use
\begin{equation}
 v_j^{\mathrm{src}}
 =\max\bigl\{v_j^{\mathrm{src,bulk}},\ 
                  \sup_T v_j^{\mathrm{src},T}\bigr\}.
 \label{eq:offset-family-source-discrepancy}
\end{equation}
For a bulk-only comparison the supremum is omitted.  Otherwise it
ranges over the periods being compared, or a compatible admitted
family, with the same source weights and caps.  This simultaneous
discrepancy is uniformly bounded by the caps.  It always includes the
bulk data that determine $c_j$; a discrepancy on an isolated torus
is not used to control that scalar.

\begin{proposition}[Source transport along offset trajectories]
\label{prop:source-input}
Every admitted run above carries the microscopic source
$\exp(\sum_x z_x\cdot q_x)$ exactly through its block integrations.
The additional source lists begin at zero and satisfy their caps at
each layer.  A step of side $l_j$, from layer $j$ to layer $j-1$, has a
bulk coefficient $c_j>0$ with
\begin{equation}
 z_x=\frac{z'_{[x]}}{l_j^2c_j},\qquad
 |c_j-1|\le C_LR_j^{\mathrm{src}}.
 \label{eq:offset-source-substitution}
\end{equation}
For a comparison on a common ordinary step,
\begin{align}
 v_{j-1}^{\mathrm{src}}
 &\le q_s v_j^{\mathrm{src}}
 +C_L(\log H_j)^D
        (u_j+\epsilon_h+|\lambda_j|/H_j),\notag\\
 |c_{2,j}-c_{1,j}|
 &\le C_LR_j^{\mathrm{src}}
 \bigl[v_j^{\mathrm{src}}+(\log H_j)^D
        (u_j+\epsilon_h+|\lambda_j|/H_j)\bigr],
 \qquad q_s<1.
 \label{eq:offset-source-comparison}
\end{align}
The statements hold in bulk and on every compatible admitted period.
The coefficient $c_j$ is the same bulk coefficient in all these periods.
\end{proposition}

\begin{proof}
The initial source has zero additional regular and covering lists and
has all support, normalization, reality, and internal covariance
properties required in Section~\ref{sec:sources}.
Proposition~\ref{prop:admission} supplies the coupling bands for every
step, including the inclined first step.  Hence
Proposition~\ref{prop:source-step} applies inductively and renews the
caps.  Equation~\eqref{eq:offset-source-substitution} is
Equation~\eqref{eq:orig-28}.  We spell out the simultaneous use of
the comparison proof, since the scalar is extracted in bulk and then
used on every period.

Before the final scalar division, the one-difference estimates in
the proof of Proposition~\ref{prop:source-step} give an isolated
degree-one transfer factor $C/L+o_H(1)$ and a higher-degree factor
$C/L^2+o_H(1)$; these are the discrepancy versions of
Equations~\eqref{eq:orig-32} and~\eqref{eq:orig-31}.
The remaining regular terms and covering transfers have the smaller
allowances specified there.  Apply these bounds to the bulk and to
each period, using the dominating discrepancies $u_j$ and
$v_j^{\mathrm{src}}$.  Extract $c_{2,j}-c_{1,j}$ from the bulk
degree-one output coefficients.  Its source-discrepancy contribution
retains that degree-one transfer gain; its remaining contribution is
bounded by
$C_LR_j^{\mathrm{src}}(\log H_j)^D
(u_j+\epsilon_h+|\lambda_j|/H_j)$.

Now perform the alignment subtraction and the scalar division on
each period with these bulk coefficients.  The extraction paragraph
of the source proof controls the subtraction, off-mask corrections,
and winding compensations by fixed multiples of their coefficient
bounds, with spare support exponents paying for the declared winding
records.  Changing the scalar argument in an already small
positive-degree remainder has the additional source cap and is
controlled by the spare analytic radius.  Thus these operations
preserve the preceding contraction gains up to multipliers fixed
before the choice of $L$, using the already reserved support exponents.
As in the closure of that proof, choose $L$ sufficiently large for
these multipliers and then $H$ sufficiently large for the remaining
allowances.  The resulting common contraction is still $q_s<1$.
The precise-coupling discrepancy and history forcing are shared
across the family.  Taking its supremum after these normalization
estimates proves the first line of
\eqref{eq:offset-source-comparison}; the bulk extraction estimate
proves the second line.  This is the simultaneous interpretation of
Equation~\eqref{eq:orig-29}, rather than an inference from a
finite-period discrepancy alone.
The source-free output and the source-independent bulk scalar are
unchanged.  No terminal-coupling equality is a hypothesis of that
proposition.
\end{proof}

For each run define the positive field factor
\begin{equation}
 B_{P,N,s}=\prod_{j=1}^N c_{P,N,s,j}^{-1}.
 \label{eq:offset-field-factor}
\end{equation}
We suppress some subscripts when a run has already been specified.

\begin{lemma}[Fixed-layer sources and comparison of products]
\label{lem:source-comparison}
Under \eqref{eq:field-offset-matching}, the source discrepancy tends to
zero at each fixed layer, and $c_{2,j}-c_{1,j}\to0$ at every fixed step
$j\ge1$.  In addition, for the ordinary and either inclined run at the
same $N$ and $\beta=\vartheta_N+s$, with $|s|\le R$, one has
\begin{equation}
 C_H^{-1}\le \frac{B_{\pm,N,s}}{B_{r,N,s}}\le C_H,
 \label{eq:offset-product-ratio}
\end{equation}
uniformly in $N$.  The products for the two mirror inclinations agree.
\end{lemma}

\begin{proof}
The normalized discrepancies are uniformly bounded by the caps.
For a fixed $j$, Proposition~\ref{prop:matching} makes the source-free
forcing in \eqref{eq:offset-source-comparison} tend to zero; the shared
ordinary history also gives $\epsilon_h\to0$.  Therefore the upper
limits $V_j=\limsup v_j^{\mathrm{src}}$ satisfy
$V_{j-1}\le q_sV_j$.  For every $m$, $V_j\le q_s^mV_{j+m}$;
boundedness and $q_s<1$ imply $V_j=0$.  The second line of
\eqref{eq:offset-source-comparison} proves the scalar assertion.

For the product estimate, compare the ordinary and inclined runs after
their first steps, identifying their output coordinates.  Admission
gives $|\lambda_j|\le C_{L,R'}$.  If $x=N-1-j$, the common-tail history
forcing is bounded by $C_LH_N^D r_0^x$, with $r_0<1$.  Iteration of the
shape recurrence used in Proposition~\ref{prop:matching} gives, after
enlarging a fixed exponent $D$ and choosing $r_1\in(0,1)$,
\begin{equation}
 u_j\le \min\{C,C_LH_N^D r_1^x\}
             +C_L\frac{(\log H_j)^D}{H_j}.
 \label{eq:offset-clipped-comparison}
\end{equation}
Here one first sums the geometric convolution and then uses the uniform
cap.  The logarithmic ratio is eventually decreasing as its argument
increases, so its convolution has the displayed bound.  Finitely many
smaller arguments are absorbed in $C_L$.
Applying \eqref{eq:offset-source-comparison} gives the same bound for
$v_j^{\mathrm{src}}$ and
$|c_{\pm,j}-c_{r,j}|/R_j^{\mathrm{src}}$, with possibly larger $D$ and
$r_1<1$.  The latter quantity can be clipped because each coefficient
satisfies \eqref{eq:offset-source-substitution}.

The second term is summable after multiplication by the source cap:
\[
 \sum_{j\ge1}R_j^{\mathrm{src}}
               \frac{(\log H_j)^D}{H_j}
 =\sum_{j\ge1}\frac{(\log H_j)^D}{H_j^{11/10}}<\infty.
\]
For the first term, split the steps at
$J_N=\lceil K\log H_N\rceil$.  On the first $J_N$ steps from the top,
the clipped estimate costs at most
$C_H H_N^{-1/10}\log H_N$ for large $N$.  The remaining steps cost at
most
\[
 C_H H_N^D\sum_{x\ge J_N}r_1^x
 \le C_H H_N^{D+K\log r_1}.
\]
Choose $K$ so that the exponent is negative.  The first step itself is
controlled by the individual scalar bounds, and finitely many small
depths have a uniform bound.  Thus
$\sum_{j=1}^N|c_{\pm,j}-c_{r,j}|\le C_H$.
All coefficients lie in a fixed positive neighborhood of one, so the
sum of the logarithmic differences is bounded as well.  Exponentiation
proves \eqref{eq:offset-product-ratio}.  This is the argument of
Proposition~\ref{prop:source-products}, with bounded-offset admission
supplying its bound on $\lambda_j$.  Reflection covariance of the
bulk source rule gives equality for the mirror inclinations.
\end{proof}

\subsection{Physical embeddings and moment bounds}

We give the terminal lattice spacing the physical value one.  Set
\begin{equation}
 a_{r,N}=L^{-N},\quad U_r=\mathrm{Id},\qquad
 a_{\pm,N}=\tfrac15L^{-N},\quad U_\pm=O_\pm^{-1}.
 \label{eq:offset-embeddings}
\end{equation}
Thus an inclined first step brings the coarse frame to the standard
frame.  With $a=a_{P,N}$, $U=U_P$, and $B=B_{P,N,s}$, define
\begin{equation}
 \phi_{P,N,s}(f)=Ba^2\sum_{x\in\Z^2} f(aUx)\cdot q_x.
 \label{eq:offset-fields}
\end{equation}
On a torus the sum is over its periodic microscopic lattice.  For
components $\boldsymbol\alpha=(\alpha_1,\ldots,\alpha_n)$, let
\begin{equation}
 S_{n,P,N,s}^{\boldsymbol\alpha,\mathrm{cut}}
 =(Ba^2)^n\sum_{x_1,\ldots,x_n}
       \E\!\left[\prod_{i=1}^nq_{x_i}^{\alpha_i}\right]
       \delta_{(aUx_1,\ldots,aUx_n)}.
 \label{eq:offset-moment-measures}
\end{equation}
The expectation is taken in the periodic plane state unless a torus is
specified.  We use the common physical square tori of side
$M_k=2^kM_0$ supplied by Lemma~\ref{lem:inclined-periods}.

\begin{lemma}[Uniform moment bounds]
\label{lem:moment-bound}
For $N\ge K_0$, $|s|\le R$, and every type $P$, the component moment
measures in the plane are nonnegative and satisfy
\begin{equation}
 S_{n,P,N,s}^{\boldsymbol\alpha,\mathrm{cut}}
      (Q_1\times\cdots\times Q_n)\le C^n n!
 \label{eq:offset-unit-box-bound}
\end{equation}
for arbitrary physical unit boxes.  The ordinary torus estimates are
uniform in $k$; on each fixed inclined physical torus the same assertion
holds with a constant depending on that torus.

Every fixed finite list of real bounded compactly supported vector
tests consequently has a joint exponential moment in a neighborhood
of zero, uniformly in these cutoffs.  The same is true of continuous
periodic tests on each fixed torus.
\end{lemma}

\begin{proof}
Component positivity is Lemma~\ref{lem:finite-beta}; its finite-volume
expansion also applies to the compatible oblique periods.
For an ordinary run, choose at most $J$ terminal sites and denote their
descendant-cell vector field sums by $X_Y$.  Iterating
\eqref{eq:offset-source-substitution} identifies $z_Y$ with the source
coupled to $X_Y$.  After removal of its source-independent scalar, the
terminal source density is
\[
 e^{X(V)}e^{\sum_Yz_Y\cdot V_Y+G(V;z)}\Xi_z(V).
\]
Let $Z$ be its normalizer at $z=0$, and let $\mu_0$ be that source-free
terminal probability law.  At $|z_Y^\alpha|\le1$, source anchoring gives
$\sup|G|\le JR_0^{\mathrm{src}}$.  Every changed covering label has
support meeting one of the selected sites, whence
\[
 \sum_\lambda e^{2s_\lambda}
       \norm{k_\lambda(z)-k_\lambda(0)}_\infty\le Jw_0.
\]
The two activity lists obey the combined cap with a fixed multiplier.
The base law is an ordinary positive endpoint law, reflection positive
at the block seams, with $b_0\in[H/2,2H]$.  Its periods have all the
required divisibilities.  Proposition~\ref{prop:integrated-comparison},
with the band verification in Lemma~\ref{lem:band-endpoints},
therefore gives
$Z^{-1}\int e^X|\Xi_z-\Xi_0|\le e^{Jw_0}-1$.
As in the proof of Proposition~\ref{prop:moment-majorants}, subtraction
of the endpoint-spin generating function yields the explicit bound
\begin{align}
 \left|\E e^{\sum_Yz_Y\cdot X_Y}
       -\E_{\mu_0}e^{\sum_Yz_Y\cdot V_Y}\right|
 &\le e^{3J}(e^{JR_0^{\mathrm{src}}}-1)
       +e^{3J+JR_0^{\mathrm{src}}}(e^{Jw_0}-1)\notag\\
 &=o_H(1)
 \label{eq:offset-local-source-bound}
\end{align}
for fixed $J$, uniformly in depth, offset, and ordinary torus size.

For one real component cell sum $X$, the bounds at sources $1$ and $-1$
give $\E e^{|X|}\le C_0$ and $\E|X|^n\le C_0n!$.  H\"older's
inequality bounds a product of $n$ such sums by $C^nn!$.  Every unit box
is covered by a bounded number of terminal cells.  Positivity of the
component measures extends this estimate to the product of any $n$
unit boxes, at an exponential cost in $n$.  Periodic coverings handle
boxes crossing a period.  At fixed cutoff the cell sums are bounded
local observables, so the estimates pass to the plane.

For an inclined field in the plane, use the ordinary field at the same
bare coupling and depth.  For either sign, the exact relation is
\[
 \phi_{\pm,N,s}(f)
 =\frac{B_{\pm,N,s}}{25B_{r,N,s}}
       \phi_{r,N,s}\bigl(f(O_\pm^{-1}\,\cdot/5)\bigr).
\]
Lemma~\ref{lem:source-comparison} and a bounded covering of the dilated,
rotated boxes prove \eqref{eq:offset-unit-box-bound} in the plane.

On a fixed inclined physical torus, the terminal volume is fixed.
Absolute coefficient norms and support-hit bounds bound the entire
source numerator in that volume.  The source-free normalizer is
uniformly bounded below after its bulk scalar has been removed:
integrate all spins in one sufficiently small spherical cap.  There
are then no bad bonds, the mandatory gas has empty inventory and
$\Xi=1$, and the regular and kinetic exponents are bounded.  Positivity
of the complete source-free density permits this restriction of the
integral.  The resulting bound can depend on the fixed terminal volume
and on $H$, but not on $N$ or $s$.  The inclined descendants of a
terminal site remain within a fixed distance of its center, and every
physical unit box meets descendants of only boundedly many sites.
The preceding one-cell and covering argument now proves the asserted
fixed-torus estimate.  Reflection positivity of inclined steps is not
used.

Finally consider a real linear combination $Y$ of a finite list of
field smearings.  For even $n$, expand components and integrate the
absolute test coefficients against the nonnegative measures in
\eqref{eq:offset-unit-box-bound}.  A fixed bounded support can be
covered by finitely many unit boxes, so
$\E|Y|^n\le C_Y^n n!$.  Odd absolute moments follow by Cauchy--Schwarz
from the adjacent even moments, since
$\sqrt{(n-1)!(n+1)!}\le\sqrt2\,n!$ for odd $n\ge1$.
Increasing $C_Y$ gives the estimate at every order.  The exponential
series converges for a sufficiently small positive argument.
H\"older's inequality gives the same conclusion for the sum of the
absolute values of the finite list.  This proves the assertions for
signed as well as nonnegative tests.
\end{proof}

\subsection{Convergence at a limiting offset}

The bounds just established allow us to pass from descendants of fixed
coarse cells to arbitrary compact tests.  We first work on a fixed
physical torus, where the density and source comparisons control
normalized integrals directly.  The uniform volume estimate then
identifies the plane limit.

\begin{proposition}[Limits and comparison at fixed offsets]
\label{prop:offset-limits}
Let $N\to\infty$ and $s\to s_*$ with $|s|\le R$, for a fixed type
$P$.  The fields \eqref{eq:offset-fields} converge jointly in law and
in all joint moments for every finite list of continuous compactly
supported vector tests in the plane.  The analogous assertions hold
for continuous periodic tests on each fixed common physical torus.
The limiting laws are determined by their moments.

Two such sequences satisfying \eqref{eq:field-offset-matching} have
the same limits in their specified embeddings.  At every order and
component choice, the plane measures
\eqref{eq:offset-moment-measures} converge on continuous compact tests
and in the strong topology of $\mathcal S'((\R^2)^n)$.

Denote the ordinary limiting hierarchy at offset $s\in I$ by $F_s$,
including its moment-determined joint laws.  For either inclined type,
the limit at offset $s$ is
\begin{equation}
 F_{s+d},\qquad s,s+d\in I,\qquad d=A_+-A_r=A_--A_r.
 \label{eq:inclined-offset-limit}
\end{equation}
\end{proposition}

\begin{proof}
Fix a common physical torus and a layer $j$ below the initial inclined
step, if present.  Its grid has mesh $L^{-j}$.  Align the grids across
the runs and use the same ordinary coordinate prescriptions on their
common tail.  These alignments are allowed by the translation covariance
of the blocking.  Translating a microscopic origin has, by microscopic
translation invariance, the same effect on the moment measures as a
translation of size $O(a)$.  Lemma~\ref{lem:moment-bound} and uniform
continuity make this residual error vanish on continuous tests.

Assign a test's value at a layer-$j$ center to every microscopic
descendant of that center.  The exact source multiplier for these cell
tests is
\begin{equation}
 L^{-2j}\prod_{i=1}^j c_i^{-1}.
 \label{eq:offset-cell-multiplier}
\end{equation}
Indeed its microscopic coefficient is $Ba^2$, and each traversed step
multiplies it by $l_i^2c_i$.  For the traversed steps in either geometry,
\[
 a^2\prod_{i=j+1}^N l_i^2=L^{-2j},\qquad
 B\prod_{i=j+1}^N c_i=\prod_{i=1}^j c_i^{-1}.
\]
Their product is \eqref{eq:offset-cell-multiplier}; the area factor of
the inclined first step is already included in $a^2$.
Only a fixed number of scalar factors occur, so this multiplier is
bounded, bounded away from zero, and compares between matched runs by
Lemma~\ref{lem:source-comparison}.

At this fixed layer the period and its number of sites are fixed.
The coefficient norms for canonical records sum their path
coefficients; the masked regular norms bound the sums of their
suprema; and the covering norms bound the activity sums at a support
hit.  Proposition~\ref{prop:matching} therefore makes the differences
of the regular exponents and the absolutely convergent covering series
tend to zero.  The free-history bounds give the same assertion for
the kinetic kernels.  Their masks and reference thresholds are common,
and $b_j$ remains in a fixed band.  Thus the source-free exponents
are bounded and compare uniformly after their volume scalars have
been discarded.  The source-list assertion in
Lemma~\ref{lem:source-comparison} gives the corresponding comparison
of the positive-degree terms in absolute coefficient norms.

For a finite list of cell tests choose one sufficiently small complex
polydisc in their source variables so that all substituted arguments
lie in a smaller source polydisc at layer $j$.  The absolute series
bounds control perturbations of these arguments as the multiplier
\eqref{eq:offset-cell-multiplier} varies.  The source-free denominators
have a uniform positive lower bound by the same small-cap argument
as in Lemma~\ref{lem:moment-bound}, now at layer $j$.  Since the removed
volume scalars are source independent, they cancel in the ratios.
The normalized generating functions of the cell tests consequently
differ by a quantity tending uniformly to zero on this polydisc.
Cauchy's formula proves comparison of all fixed joint moments.

Every microscopic descendant lies within $CL^{-j}$ of its layer-$j$
center, for both cell geometries in the embeddings
\eqref{eq:offset-embeddings}.  On the fixed torus, replacing a continuous
test by these cell values therefore has an error bounded by its modulus
of continuity at $CL^{-j}$.  Telescoping a product of tests one factor
at a time and using \eqref{eq:offset-unit-box-bound} bounds the moment
error by a constant times these moduli, uniformly in depth.  First let
the depths tend to infinity at fixed $j$, and then let $j\to\infty$.
This proves moment comparison for continuous tests.  Comparing two
arbitrary subsequences of one sequence proves the Cauchy criterion,
hence moment convergence.

The factorial bounds of Lemma~\ref{lem:moment-bound} imply tightness for
each finite list of real field variables and uniform integrability of
every polynomial in them.  Each subsequential law has the moments just
constructed and an exponential moment in a neighborhood of the origin.
Its moment generating function is analytic there and is determined by
those moments; this determines the law.  All subsequential laws thus
coincide, proving joint convergence in law on the fixed torus.

For the plane, let
$\exp(i\sum_{\nu=1}^r t_\nu\phi_{P,N,s}(f_\nu))$ be a characteristic
function observable.  Its modulus is one regardless of $B$, and its
microscopic support is the support of the smearings.  For sufficiently
large $k$ this support satisfies the lift and half-size conditions in
Proposition~\ref{prop:uniform-volume}, including the shifted-period
description in Lemma~\ref{lem:inclined-periods}.  That proposition gives
\begin{equation}
 \left|\E_{\mathrm{plane}}e^{i\sum_\nu t_\nu\phi(f_\nu)}
 -\E_{\mathrm{torus}(k)}e^{i\sum_\nu t_\nu\phi(f_\nu)}\right|
 \le Ce^{-c2^k},
 \label{eq:offset-characteristic-comparison}
\end{equation}
uniformly in the cutoff and offset.  Plane tightness follows from the
plane moment bound.  Every subsequential plane characteristic function
is, for each fixed $k$, within $Ce^{-c2^k}$ of the already constructed
fixed-torus limit.  Letting $k\to\infty$ identifies all subsequential
plane laws and gives their equality for matched runs.  Plane uniform
integrability then gives convergence of all joint moments.  Only bounded
observables enter \eqref{eq:offset-characteristic-comparison}; no bound
uniform in increasing inclined torus size is needed.

Finite sums of product tests are uniformly dense on a product of
compact insertion-coordinate sets.  The component measures have
uniformly bounded mass there by \eqref{eq:offset-unit-box-bound}.
The joint moment convergence therefore gives convergence against every
continuous compact test to a nonnegative Radon measure.  Summing unit
boxes with the integrable weights $(1+|y_i|^2)^{-2}$ gives, at fixed
order $n$, the common bound
\begin{equation}
 |S_n^{\boldsymbol\alpha,\mathrm{cut}}(f)|,
 \ |F_s^{\boldsymbol\alpha}(f)|
 \le C^n n!\,p_n(f),\qquad
 p_n(f)=\sup_{y_1,\ldots,y_n}|f(y)|
                    \prod_{i=1}^n(1+|y_i|^2)^2.
 \label{eq:offset-schwartz-bound}
\end{equation}
This proves temperedness and convergence on Schwartz tests by cutting
off the tails.

For completeness, this convergence is strong: it is uniform on each
bounded subset $\mathcal B$ of Schwartz space.  The higher weighted
seminorms bounded on $\mathcal B$ make the discarded tails tend to zero
uniformly in $p_n$.  On a fixed compact set, restrictions of
$\mathcal B$ are uniformly bounded and equicontinuous, hence have
finite uniform-norm nets.  Convergence for the finitely many net
elements and the common compact mass bound make the error uniform on
$\mathcal B$.  Combining these two statements proves the strong
topology assertion.

Finally, a type $\pm$ sequence with offset tending to $s$ and an
ordinary sequence with offset tending to $s+d$ satisfy
\eqref{eq:field-offset-matching}.  This proves
\eqref{eq:inclined-offset-limit}.
\end{proof}

\subsection{The relative two-point measure}

We next pass from local convergence to susceptibility and second-moment
length.  Write $S_2^{\mathrm{cut}}$ for the first-component two-point
measure of a run.  In coordinates $(y_1,z)=(y_1,y_2-y_1)$, microscopic
translation invariance gives the exact factorization
\begin{equation}
 S_2^{\mathrm{cut}}=\alpha_a\otimes\nu^{\mathrm{cut}},\qquad
 \alpha_a=a^2\sum_{x\in\Z^2}\delta_{aUx},\qquad
 \nu^{\mathrm{cut}}=B^2a^2\sum_{x\in\Z^2}
                          C_\beta(x)\delta_{aUx}.
 \label{eq:relative-factorization}
\end{equation}
The first factor tends vaguely to Lebesgue measure.  The second carries
the susceptibility and length, so we must control its mass at large
relative separation.

\begin{lemma}[Uniform two-point tails]
\label{lem:two-point-tails}
There are constants $C,c>0$, uniform over $N\ge K_0$, $|s|\le R$, and
the three types, such that for physical unit boxes with center
separation $T$,
\begin{equation}
 S_2^{\mathrm{cut}}(Q\times Q')\le Ce^{-cT}.
 \label{eq:offset-two-point-decay}
\end{equation}
Consequently, for a unit box $Q_u$ centered at $u$,
\begin{equation}
 \nu^{\mathrm{cut}}(Q_u)\le Ce^{-c|u|}.
 \label{eq:relative-unit-tail}
\end{equation}
Along each convergent-offset sequence of
Proposition~\ref{prop:offset-limits}, the measures
$\nu^{\mathrm{cut}}$ converge vaguely, in total mass, and in second
moment to a finite nonnegative measure $\nu$.  Its tensor product
with Lebesgue measure is the limiting first-component two-point
measure in the relative coordinates.
\end{lemma}

\begin{proof}
Let $X_Q$ and $X_{Q'}$ be the first-component field sums over the two
boxes.  Internal symmetry gives zero means; their fourth moments are
uniformly bounded by Lemma~\ref{lem:moment-bound}.  Write
$T_D(x)=\max\{-D,\min\{x,D\}\}$ for odd truncation.  Then
\[
 \norm{X_Q-T_D(X_Q)}_{L^2}^2
 \le D^{-2}\E|X_Q|^4\le C D^{-2}.
\]
The second moments of the other factors are bounded, so Cauchy--Schwarz
shows that replacing both variables by their truncations changes their
covariance by at most $C/D$.

For sufficiently large $T$, at least one microscopic axial direction
separates the supports by $(c_1T-C_1)/a$ lattice steps.  Apply the slab
covariance bound in Proposition~\ref{prop:uniform-volume} to the
truncated variables.  Since $aL^N$ equals either $1$ or $1/5$, it gives
\[
 |\operatorname{Cov}(T_D(X_Q),T_D(X_{Q'}))|
       \le CD^2e^{-c_2T}.
\]
Choosing $D=e^{c_2T/4}$ proves
\eqref{eq:offset-two-point-decay}; a larger constant covers bounded
separations.

Fix a unit box $Q$ near the origin.  For small $a$, $\alpha_a(Q)$ is
bounded below by a positive constant.  Positivity and
\eqref{eq:relative-factorization} imply
\[
 \alpha_a(Q)\nu^{\mathrm{cut}}(Q_u)
       \le S_2^{\mathrm{cut}}\bigl(Q\times(Q+Q_u)\bigr).
\]
The Minkowski sum on the right has a bounded unit-box covering, whose
centers lie at bounded distance from $u$.  Thus
\eqref{eq:offset-two-point-decay} proves
\eqref{eq:relative-unit-tail}.  Enlarge $K_0$ if necessary to have this
uniform lower bound for $\alpha_a(Q)$ throughout.

To identify the vague limit, choose a continuous compactly supported
function $h$ with $\int h=1$.  For every continuous compactly supported
$g$ of the relative coordinate,
\[
 S_2^{\mathrm{cut}}\bigl(h(y_1)g(y_2-y_1)\bigr)
       =\alpha_a(h)\nu^{\mathrm{cut}}(g),\qquad
 \alpha_a(h)\longrightarrow1.
\]
Proposition~\ref{prop:offset-limits} supplies the limit of the left
side.  The local bounds in \eqref{eq:relative-unit-tail} then identify
a unique nonnegative Radon measure $\nu$.  Product tests in relative
coordinates identify the limiting two-point measure with
$\dd y_1\otimes\nu$.  Finally \eqref{eq:relative-unit-tail}, summed
over boxes, makes both $1$ and $|z|^2$ uniformly integrable outside
large balls.  Vague convergence with this tail control proves
convergence of mass and second moment.
\end{proof}

For the ordinary hierarchy $F_s$, denote this relative measure by
$\nu_s$.  Its candidate susceptibility and length are
\begin{equation}
 m(s)=\nu_s(\R^2),\qquad
 \ell(s)^2=\frac{1}{4m(s)}\int |z|^2\,\dd\nu_s(z),
 \label{eq:offset-mass-length}
\end{equation}
where the second definition is made after the positivity established
below.  Finiteness already follows from Lemma~\ref{lem:two-point-tails}.

\subsection{Positive susceptibility and length}

\begin{proposition}[Canonical normalization at a fixed offset]
\label{prop:normalization}
For every $s\in I$, the quantities $m(s)$ and $\ell(s)$ in
\eqref{eq:offset-mass-length} are finite and strictly positive.
For an ordinary sequence with $N\to\infty$ and offset tending to $s$,
\begin{equation}
 B^2a^2\chi_\beta\longrightarrow m(s),\qquad
 a\xi_\beta\longrightarrow\ell(s).
 \label{eq:offset-cutoff-normalization}
\end{equation}
For an inclined sequence with limiting offset $s$ and $s,s+d\in I$,
the limits are $m(s+d)$ and $\ell(s+d)$.

Define $G_s$ by applying to $F_s$ the change of smeared fields
\begin{equation}
 \phi(f)\longmapsto\frac{1}{\ell(s)\sqrt{m(s)}}
                     \phi\bigl(f(\,\cdot/\ell(s))\bigr).
 \label{eq:canonical-hierarchy-normalization}
\end{equation}
This is the unique positive constant length and field rescaling that
sets both susceptibility and second-moment length to one.
Along the ordinary sequences in
\eqref{eq:offset-cutoff-normalization}, the canonical cutoff fields
$\Psi_\beta$ converge to $G_s$ in joint law and all joint moments, and
their component Schwinger measures converge strongly in $\mathcal S'$
at every order.  The analogous normalized inclined fields converge to
$G_{s+d}$.

The functions $m,\ell$ are continuous on $I$.  At every order and
component choice, evaluation of $F_s$ and of $G_s$ on any fixed
Schwartz test is continuous in $s$.
\end{proposition}

\begin{proof}
Choose two ordinary terminal cells in a fixed bounded region, with
positive distance between their closures.  Their two terminal sites
can be joined by a path of fixed length $r$.  On an admitted axis torus,
Corollary~\ref{cor:terminal-alignment} applies to the source-free
endpoint law in the full band $[H/2,2H]$ by
Lemma~\ref{lem:band-endpoints}.  Its hypotheses are precisely
the ordinary positive, seam-reflection-positive endpoint hypotheses
used above for \eqref{eq:offset-local-source-bound}.  Since
$t_0\to0$ and $Ht_0^2=p_0^2\to\infty$ as $H\to\infty$,
\[
 \E_{\mu_0}|V_{Y_1}-V_{Y_2}|^2
 \le r^2\bigl(t_0^2+4e^{-cHt_0^2}\bigr)=o_H(1).
\]
Internal $O(3)$ invariance gives
\[
 \E_{\mu_0}[V_{Y_1}^1V_{Y_2}^1]
   =\frac13\E_{\mu_0}[V_{Y_1}\cdot V_{Y_2}]
   =\frac13+o_H(1).
\]
Cauchy's formula applied to
\eqref{eq:offset-local-source-bound} transfers this estimate to the
cutoff cell sums $X_1,X_2$:
\begin{equation}
 \E[X_1X_2]=\frac13+o_H(1),
 \label{eq:offset-separated-cell-lower}
\end{equation}
uniformly in depth, offset, and admitted axis torus.  The large fixed
choice of $H$ makes the right side at least a constant $c_0>0$.
Passing to the plane at each fixed cutoff preserves this bound.

Choose nonnegative continuous compact tests that dominate the indicators
of the two cells and whose supports still have positive separation.
The cell descendants lie in fixed such boxes for all sufficiently large
depths.  Positivity transfers \eqref{eq:offset-separated-cell-lower}
to these continuous tests, and compact-test convergence passes it to
the limiting two-point measure.  In relative coordinates this gives
nonzero $\nu_s$-mass a positive distance from the origin.  Hence both
$m(s)>0$ and $\int|z|^2\,\dd\nu_s(z)>0$.  This also explains why a
lower bound only at coincident points would not suffice for the length.

At the cutoff, \eqref{eq:relative-factorization} has total relative
mass and length
\begin{equation}
 m^{\mathrm{cut}}=B^2a^2\chi_\beta,\qquad
 \ell^{\mathrm{cut}}=a\xi_\beta.
 \label{eq:cutoff-mass-length-identity}
\end{equation}
Lemma~\ref{lem:two-point-tails} gives convergence of their numerators
and denominators, with the latter limit positive.  This proves
\eqref{eq:offset-cutoff-normalization}.  The inclined conclusion follows
from \eqref{eq:inclined-offset-limit} and the same tail lemma.

To check normalization and uniqueness, let a translation-invariant
hierarchy have relative two-point measure $\nu$, susceptibility $m>0$,
and length $\ell>0$.  For positive constants $b,l$, the change
$\phi'(f)=b\phi(f(\cdot/l))$ sends its relative measure to
\[
 \nu'=b^2l^2(D_{1/l})_*\nu,\qquad D_{1/l}(z)=z/l.
\]
The factor $l^2$ is the Jacobian of the first insertion coordinate.
Thus $m'=b^2l^2m$ and $\ell'=\ell/l$.  Requiring $m'=\ell'=1$
forces $l=\ell$ and $b=(\ell\sqrt m)^{-1}$, which proves
\eqref{eq:canonical-hierarchy-normalization} and uniqueness among
positive constants.  An overall field sign, if permitted, gives the
same hierarchy by internal symmetry.

Apply this calculation before taking the limit, using
\eqref{eq:cutoff-mass-length-identity}.  The normalized field is exactly
\begin{equation}
 \frac{1}{\ell^{\mathrm{cut}}\sqrt{m^{\mathrm{cut}}}}
       \phi_{P,N,s}\bigl(f(\,\cdot/\ell^{\mathrm{cut}})\bigr)
  =\frac{1}{\xi_\beta\sqrt{\chi_\beta}}
       \sum_x f(U_Px/\xi_\beta)\cdot q_x.
 \label{eq:canonical-cutoff-identity}
\end{equation}
For $P=r$ this is $\Psi_\beta$; for $P=\pm$ it is the corresponding
spatial rotation of the same canonical field.

The limits of $m^{\mathrm{cut}}$ and $\ell^{\mathrm{cut}}$ are positive
and finite.  The scalar multipliers in
\eqref{eq:canonical-cutoff-identity} are therefore bounded, and the
varying dilations of a continuous compact test converge uniformly
with supports in a common compact set.  The moment bound proves that
this varying-test error tends to zero.  It also gives uniform box
bounds for the normalized measures: the inverse dilation of a unit box
is covered by boundedly many old unit boxes.  Factorial moments then
give joint-law convergence and uniform integrability just as in
Proposition~\ref{prop:offset-limits}.  The common weighted bound and
finite-net argument in that proposition give strong Schwartz
convergence for the normalized measures as well.

Finally let $s_i\to s$ in $I$.  For any one of the scalar quantities
$m,\ell$ or a fixed component Schwinger evaluation of $F$ or $G$, choose
a depth $N_i\ge i$ at offset $s_i$ for which the corresponding cutoff
quantity approximates its limit at $s_i$ within $1/i$.  The convergence
statements above apply to this sequence because its offsets tend to
$s$; they identify its limit with the quantity at $s$.  This proves
all the asserted continuities without assuming uniform convergence in
the offset in advance.
\end{proof}

\begin{lemma}[Normalization of the shooting hierarchy]
\label{lem:shooting-normalization}
The shooting hierarchy constructed in
Proposition~\ref{prop:continuum-limit} has finite positive susceptibility
and finite positive second-moment length.  Applying the unique positive
rescaling \eqref{eq:canonical-hierarchy-normalization} to it is the
limit of applying canonical susceptibility and length normalization to
its cutoff fields.  These assertions also hold when its sufficiently
large auxiliary parameters are fixed independently of those used for
the offset trajectories.
\end{lemma}

\begin{proof}
Use the shooting prescription's own constants and physical mesh
$a_N=L^{-N}$.  Theorem~\ref{thm:trace-uniform} supplies its slab
covariance estimate at scale $L^N$, and
Proposition~\ref{prop:moment-majorants} supplies uniform fourth moments
and all-order factorial bounds.  The factorization
\eqref{eq:relative-factorization} is an exact lattice identity, so the
proof of Lemma~\ref{lem:two-point-tails} applies with these constants.
Proposition~\ref{prop:continuum-limit} then identifies the relative
measure and gives convergence of its mass and second moment.
Equation~\eqref{eq:terminal-cell-moments} in the proof of
Proposition~\ref{prop:non-gaussian} gives a uniform positive two-point
expectation on two separated terminal cells.  The continuous majorants
used in the proof of Proposition~\ref{prop:normalization} show that the
limiting relative measure has positive mass away from the origin.
Its susceptibility and length are therefore strictly positive.
The varying-test argument and the exact cutoff identity
\eqref{eq:canonical-cutoff-identity} now prove the normalization claim.
All these inputs concern this shooting prescription itself, so no
identification of its auxiliary parameters with the offset construction
is required.
\end{proof}
